\section{Тестирование и результаты}

\subsection{Набор тестовых проектов}

Разработанный инструмент был протестирован на синтетических тестах и на реальных открытых проектах
аппаратного обеспечения. В тестовый набор вошли следующие проекты:

\begin{itemize}
\item Процессорные ядра архитектуры RISC-V: BlackParrot [20], CVA6 [21], Ibex [22], OpenC910 [23],
PicoRV32 [24], SCR1 [25], VeeR-EH1 [26], VeeR-EL2 [27].

\item Аппаратный блок шифрования AES [28].

\item Открытая платформа корня доверия OpenTitan [29].
\end{itemize}

Для каждого из проектов подготавливался сценарий запуска, собирающий иерархию синтезируемой части
проекта. Сценарии запуска стандартизованы и включают в себя настройку параметров сборки, указание
списков файлов и определение макросов. Результаты работы инструмента воспроизводимы при повторном
запуске в рамках одной и той же версии проекта.

\subsection{Результаты обнаружения}

Анализатор продемонстрировал корректную работу на всех оцениваемых проектах. В тестовом наборе
были зафиксированы случаи всех четырёх анализируемых типов ошибок (см. Таблицу 3).

\begin{table}[h!]
    \centering
    \begin{tabular}{|c|c|c|c|c|c|}
        \hline
        Проект & Автоматы & Мёртвые переходы & Недостиж. сост. & Дедлоки & Лайвлоки \\ \hline
        AES & 4 & 0 & 0 & 0 & 0 \\ \hline
        BlackParrot & 18 & 3 & 3 & 3 & 5 \\ \hline
        OpenC910 & 1201 & 541 & 111 & 48 & 4 \\ \hline
        CVA6 & 24 & 13 & 5 & 0 & 0 \\ \hline
        Ibex & 4 & 1 & 0 & 0 & 1 \\ \hline
        OpenTitan & 592 & 71 & 58 & 95 & 14 \\ \hline
        PicoRV32 & 2 & 1 & 1 & 0 & 0 \\ \hline
        SCR1 & 26 & 0 & 0 & 0 & 0 \\ \hline
        VeeR-EH1 & 15 & 0 & 0 & 0 & 0 \\ \hline
        VeeR-EL2 & 12 & 0 & 0 & 0 & 0 \\ \hline
    \end{tabular}
    \caption{Caption}
    \label{tab:placeholder}
\end{table}

\subsection{Классификация найденных проблем}

Ручная верификация полученных результатов показала, что обнаруженные сценарии могут быть разделены на три категории:

\textbf{Конфигурационные артефакты.} Большинство мёртвых переходов и недостижимых состояний связаны с
параметризованными модулями, инстанцированными с определёнными значениями параметров. Когда параметр
отключает некоторую функциональность, соответствующие ветви в логике перехода становятся заведомо
мёртвыми. Такие срабатывания не являются ошибками в прямом смысле слова, но их обнаружение
полезно --- оно показывает, какая часть кода <<мертва>> при данной конфигурации и, соответственно, может
быть исключена при анализе или оптимизирована при синтезе.

\textbf{Намеренное терминальное поведение.} Многие из обнаруженных лайвлоков соответствуют многофазным
автоматам, в которых фаза инициализации однократно переходит в основную фазу выполнения, никогда не возвращаясь
в начальное состояние. Согласно графовому определению, это ТССК с более чем одной вершиной, и
классифицируется как лайвлок. Однако в проектировании это закономерное поведение: автомат выполняет
инициализацию один раз и дальше работает в установившемся режиме. Аналогично, некоторые дедлоки
соответствуют терминальным состояниям обработки ошибок, из которых восстановление не предусмотрено
намеренно (например, состояние «аппаратная блокировка» после обнаружения необратимого сбоя).
Инструмент помечает такие случаи для ручной проверки, позволяя разработчику подтвердить намеренность.

\textbf{Потенциальные ошибки.} Меньшая часть обнаруженных аномалий представляет действительные ошибки,
требующие исправления. Так, в проекте OpenC910 мёртвые переходы, вызванные константно-значными сигналами,
привели к выявлению фрагмента кода, который никогда не исполняется; соответствующая проблема была оформлена и
передана в основной репозиторий проекта [30].

Такая классификация характерна для результатов статического анализа в целом: инструмент не может
самостоятельно различить намеренную и непреднамеренную логику, но предоставляет разработчику
структурированное представление для быстрого принятия решений. Для облегчения этого процесса результаты
каждого анализа выводятся в формате DOT-графа с аннотациями, восстановленными из исходного кода.

\subsection{Производительность}

Для оценки вычислительных характеристик предложенного подхода было измерено сквозное время анализа
(включая разбор, трансляцию в RTLIL, извлечение автоматов и графовый анализ) на репрезентативном
подмножестве проектов. Измерения проводились на системе AMD Ryzen 5 4600G (6 ядер, 12 потоков, частота
до 4,3 ГГц) с 16 ГБ оперативной памяти. Результаты представлены в Таблице 4.

% Requires: \usepackage{array}
\begin{table}[h!]
    \centering
    \caption{Сквозное время анализа для выбранных проектов}
    \label{tab:perf1}
    \begin{tabular}{|l|c|c|c|}
        \hline
        Проект & Строк кода & Автоматы & Время выполнения \\ \hline
        PicoRV32  & 7{,}8 тыс.  & 2    & 0{,}6 c  \\ \hline
        BlackParrot & 148 тыс.   & 18   & 27 c    \\ \hline
        OpenC910   & 406 тыс.   & 1201 & 121 мин \\ \hline
    \end{tabular}
\end{table}

Время анализа изменяется в широких пределах: от долей секунды для небольших проектов (PicoRV32) до
27 секунд для проекта среднего размера (BlackParrot) и примерно двух часов для самого крупного
(OpenC910, более 1200 автоматов на 406 тыс. строк кода).

Для выявления узкого места масштабируемости в Таблице 5 приведено разбиение времени выполнения по фазам
для проекта OpenC910. Фаза графового анализа (включая обнаружение дедлоков, лайвлоков, недостижимых
состояний и восстановление имён) завершается менее чем за две минуты (2\% от общего времени) даже для
этого самого крупного теста. Основные затраты приходятся на извлечение автоматов (55\%) и их обнаружение
(34\%), что в сумме составляет 89\% общего времени. Это согласуется с обсуждавшейся в разделе 5.2
худшей сложностью процедуры извлечения.

% Requires: \usepackage{booktabs}
\begin{table}[h]
    \centering
    \caption{Разбиение времени выполнения по фазам для проекта OpenC910 (общее время 121 мин)}
    \label{tab:perf2}
    \begin{tabular}{|l|c|c|}
        \hline
        \textbf{Фаза} & \textbf{Время} & \textbf{Доля} \\
        \hline
        Извлечение автоматов (fsm\_extract) & 67 мин & 55\% \\
        Обнаружение автоматов (fsm\_detect) & 41 мин & 34\% \\
        Сопоставление с кодом (fsm\_match\_conditions) & 5 мин & 4\% \\
        Разбор и трансляция в RTLIL & 3 мин & 2\% \\
        Графовый анализ и восстановление имён (fsm\_analyze) & 2 мин & 2\% \\
        Оптимизационные проходы (opt) & 3 мин & 3\% \\
        \hline
    \end{tabular}
\end{table}

Таким образом, узким местом масштабируемости является исчерпывающий перебор состояний при извлечении
автомата, а не предложенные в работе алгоритмы анализа. Предложенная в работе сложность обеспечивает
практически приемлемые времена выполнения для интеграции в рабочие процессы верификации — в частности,
в ночные регрессионные прогоны или проверки перед коммитом.

\subsection{Обсуждение проекта OpenC910}

Проект OpenC910 заслуживает отдельного обсуждения. Анализатор обнаружил в нём 1201 экземпляр автомата, из
которых только 88 являются уникальными (соответствуют различным определениям автоматов в исходном коде).
Остальные экземпляры — это дубликаты, возникшие при элаборации, поскольку анализатор работает на уровне
netlist, где для каждого экземпляра модуля создаётся отдельный набор логических элементов. Такое
дублирование также затрагивает и сообщения об ошибках: одна и та же проблема в одном определении
автомата повторяется для всех его экземпляров.

Для сравнения: базовая реализация Yosys (версия 0.58) без наших расширений обнаружила в том же проекте
только 446 экземпляров (43 уникальных). Таким образом, предложенный инструмент выявляет примерно вдвое
больше уникальных автоматов (88 против 43), что объясняется расширениями, описанными в разделе 5.2, ---
в первую очередь, разбиением проводов на части, поддержкой автоматов Мили и обработкой синхронного
сигнала сброса. Эти расширения позволяют обнаруживать автоматы, которые немодифицированный Yosys
распознать не в состоянии.

В этом проекте было обнаружено значительное число мёртвых переходов, и соответствующие проблемы
переданы в основной репозиторий проекта. Как упоминалось в подразделе 6.3, переменная ace\_cfig,
жёстко привязанная к нулю, является характерным примером конфигурационного артефакта, приводящего
к появлению мёртвых переходов.

\subsection{Сравнение с коммерческим инструментом}

Для синтетических тестов и проекта PicoRV32 выполнено сравнение с коммерческим инструментом статического
анализа по точности извлечения автоматов и обнаружения ошибок. Результаты показали, что разработанный
инструмент совпал с коммерческим аналогом во всех стандартных случаях и, сверх того, обнаружил существенно
больше экземпляров автоматов в случаях, когда логика переходов выходит за рамки оператора case — например,
когда блок case и связанная с ним логика сброса объединены в одном блоке always.

Особенно показателен случай, обнаруженный в процессорном ядре VeeR-EL2. В этом проекте используется
характерный шаблон, когда комбинационная логика вычисления следующего состояния автомата (с помощью
оператора case внутри блока always\_comb) располагается в одном процедурном блоке, а регистр состояния
реализуется как отдельный параметризованный модуль триггера (rvdff), подключаемый через порты модуля.
Упрощённая версия этого шаблона приведена ниже.

\begin{lstlisting}
module top (input logic clk, rst_l, en);
  typedef enum logic [2:0]
    {IDLE, START_WAIT, STOP} state_t;
  state_t buf_state, buf_nxtstate;
 
  always_comb begin
    buf_nxtstate = IDLE;
    case (buf_state)
      IDLE:       buf_nxtstate = START_WAIT;
      START_WAIT: buf_nxtstate = STOP;
      STOP:       buf_nxtstate = IDLE;
    endcase
  end
 
  rvdff #(.WIDTH($bits(state_t)))
    buf_state_ff (.din(buf_nxtstate),
                  .dout(buf_state),
                  .clk(clk), .rst_l(rst_l));
endmodule
 
module rvdff #(parameter WIDTH=1)
  (input  logic [WIDTH-1:0] din,
   input  logic clk, rst_l,
   output logic [WIDTH-1:0] dout);
  always_ff @(posedge clk or negedge rst_l)
    if (!rst_l) dout <= 0;
    else        dout <= din;
endmodule
\end{lstlisting}

Этот шаблон фактически распределяет конечный автомат по границам модулей: комбинационный always\_comb
блок вычисляет следующее состояние, а параметризованный триггерный модуль хранит текущее. Коммерческий
инструмент не распознал такой шаблон как автомат, поскольку АСД-ориентированное обнаружение ожидает, что
регистр состояния и логика переходов находятся в одном и том же модуле. Предложенный в работе подход,
основанный на netlist, корректно определил автомат, так как после элаборации и трансляции в RTLIL
инстанцирование триггера уже разрешено, и мультиплексорное дерево напрямую связано с триггером ---
независимо от исходной иерархии модулей. Этот пример наглядно демонстрирует преимущества подхода,
работающего с уровнем netlist, и показывает практическую ценность расширений, реализованных в настоящей работе.

\subsection{Обработка параметризированных модулей}

Дополнительно следует отметить важность корректной обработки параметризованных автоматов. В современных
проектах на SystemVerilog часто используются параметризованные модули, в которых один и тот же фрагмент
кода может порождать автоматы с различающимся поведением в зависимости от значений параметров. Анализ на
уровне модуля без учёта конкретных параметров в общем случае недостаточен.

Поскольку разработанный инструмент работает на уровне netlist, генерируемого после элаборации параметров,
каждый экземпляр автомата анализируется независимо со своими конкретными значениями параметров. Это
позволяет обнаруживать конфигурационно-зависимое поведение, включая состояния и переходы, достижимые в
одном экземпляре, но недостижимые в другом.

Такие эффекты были зафиксированы в ходе экспериментального исследования. Например, в модуле bp\_uce [19]
проекта BlackParrot разные настройки параметров порождали разные графы автомата: в одних конфигурациях
некоторые состояния оказывались недостижимыми. Это иллюстрируется на Рисунке 3, где показаны два экземпляра
одного и того же автомата с различными наборами параметров. Кроме того, в ряде автоматов наблюдалась
многофазная работа, состоящая из фазы инициализации и основной фазы выполнения, которая никогда не
возвращается в начальное состояние; такое поведение, будучи намеренным, тем не менее классифицируется как
лайвлок по графовому определению и помечается для ручной проверки.

\begin{figure}[h!]
    \centering
    \includegraphics[width=0.7\linewidth]{black_parrot.png}
    \caption{Два экземпляра одного и того же автомата с разными значениями параметров в модуле bp\_uce
    проекта BlackParrot: (вверху) без недостижимых состояний, (внизу) с недостижимыми состояниями}
    \label{fig:black_parrot}
\end{figure}

Работа на уровне netlist, полученного после элаборации, позволяет разработанному методу улавливать
конфигурационно-зависимое поведение и выявлять сценарии, которые были бы невидимыми при анализе на уровне
исходного кода или отдельных модулей.

\section{Заключение}
\label{sec:Chapter5} \index{Chapter5}

В настоящей работе разработан инструмент SVAN-FSM --- автоматический анализатор конечных автоматов в
описаниях на языке SystemVerilog, созданный как часть системы статического анализа SVAN. Инструмент
выявляет четыре класса проблем, связанных с конечными автоматами: мёртвые переходы, недостижимые состояния,
дедлоки и лайвлоки. Дополнительно обнаруживаются взаимодействующие автоматы, что позволяет выявлять
глобальные зависимости между блоками управления. Результаты анализа представляются в виде DOT-графов,
связанных с исходным кодом.

Работая на уровне netlist, анализатор сочетает графовый анализ с сопоставлением результатов с исходным
кодом. Это обеспечивает надёжную обработку самых разных стилей написания RTL-кода, включая случаи, когда
логика автомата распределена по нескольким процедурным блокам или даже по модульным границам. В
инфраструктуру конечных автоматов Yosys внесены несколько существенных улучшений: поддержка синхронного
сигнала сброса, разбиение проводов на части, единообразная обработка автоматов Мура и Мили.

Экспериментальное исследование на 10 открытых проектах аппаратного обеспечения, включая процессорные ядра
RISC-V и критичные по безопасности проекты, продемонстрировало применимость подхода к промышленным
масштабам. Время анализа остаётся приемлемым даже для самых крупных тестовых проектов: для OpenC910 (406 тыс.
строк кода, 1201 автомат) полный анализ занимает около двух часов на стандартной рабочей станции.

Сравнение с коммерческим статическим анализатором показало, что предложенный инструмент значительно
превосходит аналог по охвату обнаружения — в первую очередь в проектах с нетиповыми шаблонами реализации,
например в случаях, когда регистр состояния вынесен в отдельный модуль.

Одним из текущих ограничений инструмента является использование эвристики распространения констант для
обнаружения мёртвых переходов. Дальнейшие работы будут направлены на интеграцию более точных методов, таких
как рассуждения на основе SAT-решения и символьное исполнение, для обнаружения мёртвых переходов,
включающих непостоянные входы.

Другие запланированные направления включают:

\begin{itemize}
\item Расширение анализа взаимодействия автоматов для обнаружения глобальных сценариев дедлоков (дедлоков
уровня системы, а не отдельного автомата).

\item Интеграцию правил кодирования STARC, охватывающих типовые рекомендации по стилю описания автоматов на
SystemVerilog.

\item Расчёт дополнительных метрик автоматов: цикломатической сложности графа, стиля кодирования состояний
(двоичный, Грея, one-hot), глубины графа, числа неиспользованных битов регистра состояния.

\item Интеграция с пользовательским интерфейсом SVAN для более удобной навигации между результатами анализа
и исходным кодом.
\end{itemize}
